\chapter{Heuristische Evaluation der Anwendung}
\label{chapter:kapitel4}
Durchführung der eigentlichen Evaluation auf Basis eines operationalisierten
Kriterienkatalogs, Dokumentation der Befunde mit Schweregradbewertung und
Ableitung von Optimierungspotenzialen.

\section{Untersuchungsdesign und Operationalisierung}
\label{section:untersuchungsdesign}

Die Untersuchung erfolgt als heuristische Evaluation in Form einer
Experteninspektion.
Statt die Anwendung durch Endnutzer erproben zu lassen, prüft der Evaluator sie
systematisch gegen einen zuvor festgelegten
Kriterienkatalog.\footcite[Vgl.][S.~254 f.]{nielsen1990heuristic}
Dazu durchläuft er die für die Bilddokumentation relevanten Aufgabenpfade
mehrfach und gleicht die einzelnen Dialogelemente mit den Prüfkriterien ab.
Der Inspektionsumfang beschränkt sich auf die in
Kapitel~\ref{section:untersuchungsgegenstand} beschriebenen Funktionen.
Dass der Verfasser die identifizierten Abweichungen im Anschluss selbst als
Entwickler behebt, ist eine methodische Einschränkung und wird in
Kapitel~\ref{grenzen} reflektiert.

Als Bewertungsmaßstab dienen die sieben Interaktionsprinzipien der ISO~9241-110
(Kapitel~\ref{grundsaetzeISO}).
Da sie abstrakt formuliert sind, werden sie für die Evaluation operationalisiert,
also in überprüfbare Prüfkriterien überführt.
Grundlage ist die normeigene Anwendungstabelle (Tabelle~A.1), die zu jedem
Prinzip konkrete Gestaltungsempfehlungen
aufführt.\footcite[Vgl.][Tabelle~A.1]{DIN_ISO_9241-110:2020}
Jede Empfehlung wird zunächst auf ihre Anwendbarkeit im vorliegenden
Nutzungskontext geprüft.
Nicht anwendbare Empfehlungen werden ausgeschlossen, wie es die Norm vorsieht.
Von den 64 Einzelempfehlungen entfallen so 22.
Die verbleibenden 42 werden zu 41 aufgabenbezogenen Prüffragen operationalisiert,
die sich mit \enquote{erfüllt} oder \enquote{Abweichung} beantworten lassen,
wobei Kriterium~7.4 zwei Empfehlungen zur Vertrauenswürdigkeit
bündelt.\footcite[Vgl.][Tabelle~A.1]{DIN_ISO_9241-110:2020}
Der Kriterienkatalog ist in Anhang~\ref{anhang:kriterienkatalog} vollständig
dokumentiert und lässt sich grundsätzlich auf vergleichbare mobile
SAP-Fiori-Anwendungen übertragen.

Jede festgestellte Abweichung wird als Usability-Problem dokumentiert, dem
verletzten Prinzip zugeordnet und nach \ac{SG} eingestuft.
Dafür dient die fünfstufige Skala nach Nielsen, die von
\enquote{kein Usability-Problem} (0) bis \enquote{Usability-Katastrophe} (4)
reicht.\footcite[Vgl.][S.~47]{nielsen1994heuristic}
Die Einstufung richtet sich nach Häufigkeit, Schwere der Auswirkung und
Persistenz des Problems.\footcite[Vgl.][S.~47]{nielsen1994heuristic}

\begin{table}[htbp]
\centering
\caption[Schweregradskala nach Nielsen]{Schweregradskala nach
Nielsen\footcite[Mit Änderungen entnommen aus:][S.~47]{nielsen1994heuristic}}
\label{tab:schweregrade}
\begin{tabularx}{0.98\textwidth}{|c|>{\RaggedRight}X|}
\hline
\textbf{Grad} & \textbf{Bedeutung} \\\hline
0 & Kein Usability-Problem \\\hline
1 & Kosmetisches Problem - Behebung nur bei freier Kapazität \\\hline
2 & Geringfügiges Problem - niedrige Behebungspriorität \\\hline
3 & Schwerwiegendes Problem - hohe Behebungspriorität \\\hline
4 & Usability-Katastrophe - Behebung vor Auslieferung zwingend \\\hline
\end{tabularx}
\end{table}

Um die Forschungsfrage nach dem \emph{Inwieweit} der Erfüllung zu beantworten,
wird je Prinzip ein Erfüllungsgrad aus den zugehörigen Prüfkriterien abgeleitet.
Ein Prinzip gilt als \emph{erfüllt}, wenn keine oder nur kosmetische Abweichungen
(Schweregrad~$\leq 1$) vorliegen, als \emph{teilweise erfüllt} bei Abweichungen
bis Schweregrad~3 und als \emph{nicht erfüllt}, sobald eine Abweichung des
Schweregrades~4 oder mehrere des Schweregrades~3 auftreten.
Diese Zuordnung macht die Bewertung in Kapitel~\ref{section:ergebnisse}
nachvollziehbar und vergleichbar.


\section{Ergebnisse und Schweregradbewertung}
\label{section:ergebnisse}

Geprüft wurden die 41 im Nutzungskontext anwendbaren Kriterien des
Kriterienkatalogs.
Bei 11 davon wurde eine Abweichung festgestellt, 30 waren erfüllt.
Die Abweichungen verteilen sich auf zwei Befunde des Schweregrades~1, fünf des
Schweregrades~2 und vier des Schweregrades~3.
Eine Usability-Katastrophe (Schweregrad~4) trat nicht auf
(vgl. Tabelle~\ref{tab:schweregrade}).
Die Darstellung erfolgt prinzipienweise.
Je Prinzip werden zunächst die erfüllten Kriterien genannt, dann die Abweichungen
einzeln, jeweils mit verletzter Normanforderung, konkreter Beobachtung und
Begründung der Schweregradeinstufung.
Abschließend wird der Erfüllungsgrad nach der in
Kapitel~\ref{section:untersuchungsdesign} festgelegten Logik bestimmt.\\
\textbf{Prinzip 1 -- Aufgabenangemessenheit:}
Für die Aufgabenangemessenheit wurden fünf Kriterien geprüft, keines wies eine
Abweichung auf.
Die Anwendung macht ihren Zweck, die fotografische Dokumentation zu einer
bestehenden Instandhaltungsmeldung, hinreichend deutlich (Kriterium~1.1), und in
jedem Schritt des Ablaufs, von der Meldungsauswahl über die Fotoaufnahme bis zur
Fotoverwaltung, stehen die erforderlichen Steuerelemente und Informationen bereit
(Kriterium~1.2).
Da der Funktionsumfang bewusst eng auf den Dokumentationszweck zugeschnitten ist
(vgl. Kapitel~\ref{section:untersuchungsgegenstand}), enthält die Anwendung
nichts, was von der Aufgabe ablenkt (Kriterium~1.3).
Sie bietet Vorbelegungen dort, wo sie zweckmäßig sind, und vermeidet solche, die
eine Fehlzuordnung der Fotos begünstigen könnten (Kriterien~1.4 und~1.5).
Das Prinzip gilt als \emph{erfüllt}.\\
\textbf{Prinzip 2 -- Selbstbeschreibungsfähigkeit:}
Von sieben Kriterien waren vier erfüllt.
Die erforderlichen Steuerelemente lassen sich zuverlässig auffinden
(Kriterium~2.2), der Bearbeitungsfortschritt ist über die Zahl erfasster Fotos
und den Bezug zur gewählten Meldung erkennbar (Kriterium~2.5), das Übernehmen
oder Verwerfen eines Fotos wird zurückgemeldet (Kriterium~2.6), und es ist
erkennbar, dass zunächst eine Meldung auszuwählen ist, bevor die Fotofunktionen
bereitstehen (Kriterium~2.7).
Drei Kriterien wiesen Abweichungen auf.
\newpage

Nach der Norm soll ein System eindeutig anzeigen, wo sich Benutzende befinden,
welche Handlungen möglich und welche Einstellungen wirksam
sind.\footcite[Vgl.][Abschnitt~5.2.2.2]{DIN_ISO_9241-110:2020}
In der Meldungssuche ist der wirksame Zustand nur eingeschränkt erkennbar.
Welche Filter aktiv sind, steht abseits im oberen rechten Fensterbereich und ist
leicht zu übersehen.
An einzelnen Bedienschritten fehlen zudem Hinweise auf die möglichen Aktionen
(Kriterium~2.1).
Da die Meldungssuche den Einstieg in jeden Vorgang bildet, tritt der Mangel
regelmäßig auf.
Seine Auswirkung bleibt moderat, weil ein übersehener Filter zu einer leeren oder
unvollständigen Trefferliste führt, den Vorgang aber nicht blockiert.
Die Abweichung wird als Schweregrad~2 eingestuft.

Interaktive Elemente sollen klar als solche erkennbar
sein.\footcite[Vgl.][Abschnitt~5.2.2.4]{DIN_ISO_9241-110:2020}
Die Buttons sind jedoch uneinheitlich gestaltet.
Einer erscheint blau gefüllt mit weißer Fettschrift, ein anderer nur blau
umrandet mit dünner blauer Schrift, sodass ihr Interaktionscharakter weniger
unmittelbar erschließbar ist (Kriterium~2.3).
Da die Schaltflächen bedienbar bleiben und die Aufgabenerledigung nicht behindert
wird, gilt der Befund als kosmetische Abweichung mit Schweregrad~1.

Informationen sollen in einem den Benutzenden vertrauten Vokabular
erscheinen.\footcite[Vgl.][Abschnitt~5.2.2.5]{DIN_ISO_9241-110:2020}
Die Meldungsart wird jedoch allein als SAP-internes Kürzel wie \enquote{M1}
angezeigt, ohne Klartext (Kriterium~2.4).
Für die handwerklich geprägte, mit SAP wenig vertraute Nutzergruppe
(vgl. Kapitel~\ref{section:untersuchungsgegenstand}) bleibt damit unklar, welche
Meldungsart sich hinter dem Kürzel verbirgt.
Das Kürzel erscheint bei jeder Meldungsauswahl und betrifft mit der
Identifikation der richtigen Meldung einen zentralen Schritt.
Da sich das Vokabular bei nur etwa zweimal wöchentlicher Nutzung kaum einprägt,
wirkt der Mangel dauerhaft und kann eine Fehlauswahl begünstigen.
Weil die angezeigten Stammdaten eine Fehlauswahl aber erkennen und korrigieren
lassen (vgl. Kriterium~6.7), bleibt die Auswirkung begrenzt und die Abweichung
wird als Schweregrad~2 eingestuft.
Mit den Kriterien~2.1 und~2.4 liegen zwei Abweichungen des Schweregrades~2 vor.
Das Prinzip gilt als \emph{teilweise erfüllt}.\\
\textbf{Prinzip 3 -- Erwartungskonformität:}
Von sieben Kriterien waren fünf erfüllt.
Die Ablaufschritte entsprechen dem Aufgabenverständnis der
Instandhaltungsmitarbeitenden (Kriterium~3.1), unmittelbare Handlungen wie das
Auslösen der Kamera werden sofort zurückgemeldet (Kriterium~3.2), Elemente
unterschiedlichen Zwecks sind erkennbar abgegrenzt (Kriterium~3.5), auf
Änderungen angeschlossener Ressourcen reagiert die Anwendung angemessen
(Kriterium~3.6), und die Inhalte werden auf der Smartphone-Anzeige passend
dargestellt (Kriterium~3.7).
Zwei Kriterien wiesen Abweichungen auf.

Bei erwartbar langen Reaktionszeiten soll eine Fortschrittsinformation ausgegeben
werden.\footcite[Vgl.][Abschnitt~5.3.2.4]{DIN_ISO_9241-110:2020}
Beim Laden der Meldungsliste oder beim Hochladen von Fotos über das Standort-WLAN
bleibt eine solche Anzeige jedoch aus, sodass nicht erkennbar ist, ob das System
noch arbeitet (Kriterium~3.3).
Das kann Benutzende dazu verleiten, eine Aktion für fehlgeschlagen zu halten und
zu wiederholen.
Da die Aufgabe erschwert, aber nicht verhindert wird, gilt die Abweichung als
Schweregrad~2.

Für die Darstellung sollen den Benutzenden vertraute sprachliche Konventionen
verwendet werden.\footcite[Vgl.][Abschnitt~5.3.3.1]{DIN_ISO_9241-110:2020}
Deutsche und englische Oberflächentexte wechseln sich jedoch ohne erkennbare
Systematik ab (Kriterium~3.4), was der Erwartung einer durchgängig
deutschsprachigen Oberfläche
(vgl. Kapitel~\ref{section:untersuchungsgegenstand}) widerspricht.
Da die Sprachmischung die gesamte Anwendung durchzieht und das Verständnis der
überwiegend nicht englischsprachig arbeitenden Nutzergruppe spürbar erschwert,
wird die Abweichung als Schweregrad~3 eingestuft.
Damit gilt das Prinzip als \emph{teilweise erfüllt}.\\
\textbf{Prinzip 4 -- Erlernbarkeit:}
Alle vier geprüften Kriterien waren erfüllt.
Die Anwendung unterstützt Benutzende dabei, die nötigen Funktionen selbstständig
zu entdecken und zielführend zu nutzen (Kriterium~4.1).
Sie lässt sich gefahrlos erkunden, da jede Aufnahme vor der Übernahme in einer
Vorschau geprüft und ohne Folgen verworfen werden kann (Kriterium~4.2).
Für das Auffinden der richtigen Meldung stehen mit Filterkriterien und
Freitextsuche alternative Wege bereit (Kriterium~4.3), und Rückmeldungen lassen
die Folgen der eigenen Aktionen erkennen (Kriterium~4.4).
Das Prinzip gilt als \emph{erfüllt}.\\
\textbf{Prinzip 5 -- Steuerbarkeit:}
Alle vier geprüften Kriterien waren erfüllt.
Benutzende können den Vorgang jederzeit unterbrechen (Kriterium~5.1) und
fortsetzen, ohne Meldungsauswahl oder erfasste Fotos zu verlieren (Kriterium~5.2).
Für die Bilderfassung stehen mit Gerätekamera und Datei-Upload zwei Wege zur
Verfügung (Kriterium~5.3), und zumindest die letzte Aktion lässt sich rückgängig
machen, indem ein Foto verworfen oder wieder entfernt wird (Kriterium~5.4).
Das Prinzip gilt als \emph{erfüllt}.\\
\textbf{Prinzip 6 -- Robustheit gegen Benutzungsfehler:}
Mit fünf Abweichungen bei zehn Kriterien wies dieses Prinzip die meisten Befunde
auf.
Erfüllt waren die Vermeidung erneuter Eingabe bereits vorhandener Informationen
(Kriterium~6.1), die Meldungsauswahl über eine Suchhilfe nach dem Prinzip der
Wiedererkennung (Kriterium~6.2), die Sicherung erfasster Fotos gegen Verlust
durch System- oder Verbindungsfehler (Kriterium~6.4), die Unterstützung beim
Erkennen einer Fehlauswahl über die angezeigten Stammdaten (Kriterium~6.7) sowie
der geringe Korrekturaufwand (Kriterium~6.10).
Die fünf Abweichungen betreffen Fehlervermeidung, Fehlererkennung und
Fehlerbehebung.

Es sollen nur aktuell gültige Auswahlmöglichkeiten zugelassen
werden.\footcite[Vgl.][Abschnitt~5.6.2.3]{DIN_ISO_9241-110:2020}
In der Meldungssuche lassen sich jedoch einander ausschließende Filter
kombinieren.
Nach Auswahl einer Kostenstelle (etwa KST6336) bleibt eine Maschine (etwa
Index G-550) wählbar, die dieser Kostenstelle nicht zugeordnet ist, sodass die
Suche zwangsläufig ohne Ergebnis bleibt (Kriterium~6.3).
Der Fall tritt nur bei bestimmten Kombinationen auf und lässt sich durch Anpassen
der Filter auflösen.
Die Abweichung wird als Schweregrad~2 eingestuft.

Eingabefehler sollen bereits vor der Weiterverarbeitung erkannt
werden.\footcite[Vgl.][Abschnitt~5.6.2.5]{DIN_ISO_9241-110:2020}
Beim Hochladen einer Datei ungeeigneten Formats bleibt jede Rückmeldung aus.
Die Datei wird ohne Hinweis verworfen (Kriterium~6.5).
Benutzende gehen dann davon aus, das Foto sei erfasst, während es tatsächlich
fehlt, sodass eine Lücke in der Dokumentation entsteht, die womöglich erst bei
der späteren Schadensanalyse auffällt.
Wegen dieser gravierenden und ohne Rückmeldung nicht erkennbaren Auswirkung wird
die Abweichung als Schweregrad~3 eingestuft.

Vor Handlungen mit schwerwiegenden Auswirkungen soll das System eine Bestätigung
einholen.\footcite[Vgl.][Abschnitt~5.6.2.6]{DIN_ISO_9241-110:2020}
Ein bereits erfasstes Foto wird jedoch mit einem einzigen Klick und ohne Rückfrage
gelöscht (Kriterium~6.6).
Da es sich vor Ort nicht ohne Weiteres erneut aufnehmen lässt, der dokumentierte
Zustand am Objekt besteht dann möglicherweise nicht mehr, kann eine Fehlbedienung
unmittelbar zu unwiederbringlichem Datenverlust führen.
Der Verlust setzt zwar eine Fehlbedienung voraus und tritt selten auf.
Wegen der Unumkehrbarkeit wird die Abweichung aber nicht als kosmetisch, sondern
als Schweregrad~2 eingestuft.

Die beiden verbleibenden Abweichungen gehen auf dasselbe Fehlverhalten bei
Fehlermeldungen zurück.
Bei manchen Fehlern erscheint zwar ein Popup, es enthält aber nur den Hinweis
\enquote{Fehlermeldung kann nicht geladen werden} statt des eigentlichen
Fehlertextes.
Damit ist die Anforderung, präzise und ohne Fachkenntnisse verständliche
Fehlermeldungen
auszugeben,\footcite[Vgl.][Abschnitt~5.6.4.2]{DIN_ISO_9241-110:2020} nicht erfüllt
(Kriterium~6.8).
Das Popup zeigt an, dass ein Problem vorliegt, benennt aber weder Art noch
Ursache, sodass Benutzende genau in der Störungssituation ohne verwertbare
Rückmeldung dastehen.
Wegen dieser in jedem Fehlerfall auftretenden Intransparenz wird der Befund als
Schweregrad~3 eingestuft.
Zugleich soll eine Fehlermeldung konstruktive Erläuterungen zu ihrer Behebung
liefern.\footcite[Vgl.][Abschnitt~5.6.4.3]{DIN_ISO_9241-110:2020}
Die Meldung nennt jedoch keine Handlungsanweisung, mit der Benutzende den Fehler
selbst beheben könnten (Kriterium~6.9).
Da sie ohne eine solche Anweisung den Vorgang nicht fortsetzen können und auf
fremde Unterstützung angewiesen sind, wird auch dieser Befund als Schweregrad~3
eingestuft.
Mit den Kriterien~6.5, 6.8 und~6.9 liegen drei Abweichungen des Schweregrades~3
vor.
Das Prinzip gilt als \emph{nicht erfüllt}.\\
\textbf{Prinzip 7 -- Benutzerbindung:}
Von vier Kriterien waren drei erfüllt.
Die Anwendung bestätigt die erfolgreiche Erfassung und Zuordnung der Fotos und
macht erkennbar, dass keine offenen Schritte verbleiben (Kriterium~7.1).
Sie benennt die Fotos automatisch aus Meldungsnummer und fortlaufender Nummer und
ordnet sie der Meldung zu (Kriterium~7.2), und die nachvollziehbare Ablage
rechtfertigt das Vertrauen in eine korrekte und zuverlässige Speicherung
(Kriterium~7.4).
Eine Abweichung wurde festgestellt.

Die Anwendung soll keine unverhältnismäßigen Anforderungen an die Benutzenden
stellen.\footcite[Vgl.][Abschnitt~5.7.2.7]{DIN_ISO_9241-110:2020}
Die Bedienung ist jedoch nicht auf die erschwerten Bedingungen in der Produktion
ausgelegt.
Insbesondere fehlt eine Anpassung an die Bedienung mit Handschuhen, die situativ
getragen werden (Kriterium~7.3,
vgl. Kapitel~\ref{section:untersuchungsgegenstand}).
Da die Handschuhe nur zeitweise getragen werden, sich die Erschwernis durch
kurzes Ablegen umgehen lässt und für die Dokumentation genügend Zeit bleibt, gilt
die Abweichung als kosmetischer Mangel mit Schweregrad~1.
Damit gilt das Prinzip als \emph{erfüllt}.

\bigskip
Tabelle~\ref{tab:erfuellungsgrad} fasst die Bewertung zusammen.
Von den sieben Prinzipien sind vier erfüllt, zwei teilweise erfüllt und eines,
die Robustheit gegen Benutzungsfehler, nicht erfüllt.
Die Abweichungen konzentrieren sich auf Selbstbeschreibungsfähigkeit,
Erwartungskonformität und vor allem Robustheit gegen Benutzungsfehler, während
die aufgabenbezogenen Prinzipien Aufgabenangemessenheit, Erlernbarkeit und
Steuerbarkeit vollständig erfüllt sind.

\begin{table}[htbp]
\centering
\caption{Erfüllungsgrad der Interaktionsprinzipien im Überblick}
\label{tab:erfuellungsgrad}
\begin{tabularx}{\textwidth}{|>{\RaggedRight}X|c|c|c|c|}
\hline
\textbf{Prinzip} & \textbf{Geprüft} & \textbf{Abw.} & \textbf{Max. SG} & \textbf{Erfüllungsgrad} \\\hline
1 Aufgabenangemessenheit & 5 & 0 & 0 & erfüllt \\\hline
2 Selbstbeschreibungsfähigkeit & 7 & 3 & 2 & teilweise erfüllt \\\hline
3 Erwartungskonformität & 7 & 2 & 3 & teilweise erfüllt \\\hline
4 Erlernbarkeit & 4 & 0 & 0 & erfüllt \\\hline
5 Steuerbarkeit & 4 & 0 & 0 & erfüllt \\\hline
6 Robustheit gegen Benutzungsfehler & 10 & 5 & 3 & nicht erfüllt \\\hline
7 Benutzerbindung & 4 & 1 & 1 & erfüllt \\\hline
\end{tabularx}
\end{table}


\section{Optimierungspotenziale}
\label{section:optimierungspotenziale}

Aus den in Kapitel~\ref{section:ergebnisse} dokumentierten Abweichungen lassen
sich konkrete Optimierungspotenziale ableiten, erfüllte Kriterien erfordern keine
Maßnahme.
Die Priorisierung folgt dem Schweregrad, der nach Nielsen unmittelbar den
Behebungsaufwand priorisiert.\footcite[Vgl.][S.~47]{nielsen1994heuristic}
Abweichungen mit gemeinsamer Ursache oder gemeinsamem Themenfeld werden zu einem
Maßnahmenbündel zusammengefasst, das die Priorität seiner schwerwiegendsten
Abweichung erhält.
Bündel mit Schweregrad~3 werden mit hoher, solche mit maximal Schweregrad~2 mit
mittlerer und rein kosmetische Befunde mit niedriger Priorität eingeordnet
(vgl. Tabelle~\ref{tab:optimierung}).

\textbf{Robustheit gegen Benutzungsfehler:}
Der größte Handlungsbedarf besteht beim einzigen nicht erfüllten Prinzip.
Beim Hochladen einer Datei ungeeigneten Formats sollte die Anwendung das Format
prüfen und unmittelbar zurückmelden, statt die Datei unbemerkt zu verwerfen
(Kriterium~6.5).
Das erkennt Eingabefehler vor der
Weiterverarbeitung\footcite[Vgl.][Abschnitt~5.6.2.5]{DIN_ISO_9241-110:2020} und
verhindert unbemerkte Lücken in der Bilddokumentation.
Die bei Ladefehlern ausgegebene Meldung sollte den eigentlichen Fehlertext
zuverlässig anzeigen, ihn zielgruppengerecht formulieren und um eine konkrete
Handlungsanweisung ergänzen, etwa den Hinweis, dass die Meldungsliste nicht
geladen werden konnte, verbunden mit der Empfehlung, den Vorgang erneut zu starten
(Kriterien~6.8 und~6.9).
So entstehen präzise, verständliche Fehlermeldungen mit konstruktiven
Erläuterungen.\footcites[Vgl.][Abschnitt~5.6.4.2]{DIN_ISO_9241-110:2020}[Vgl.][Abschnitt~5.6.4.3]{DIN_ISO_9241-110:2020}
Vor dem Löschen eines erfassten Fotos sollte eine Bestätigung abgefragt werden
(Kriterium~6.6), um unwiederbringlichen Datenverlust durch Fehlbedienung zu
vermeiden.\footcite[Vgl.][Abschnitt~5.6.2.6]{DIN_ISO_9241-110:2020}
Schließlich sollten in der Meldungssuche abhängige Filter so verknüpft werden,
dass nach einer Auswahl nur noch passende Folgewerte erscheinen (Kriterium~6.3).
Das lässt ausschließlich gültige Auswahlmöglichkeiten
zu\footcite[Vgl.][Abschnitt~5.6.2.3]{DIN_ISO_9241-110:2020} und schließt
ergebnislose Kombinationen aus.\\
\textbf{Verständlichkeit von Begriffen und Sprache:}
Ebenfalls hohe Priorität hat die Verständlichkeit für die handwerklich geprägte,
mit SAP wenig vertraute Nutzergruppe
(vgl. Kapitel~\ref{section:untersuchungsgegenstand}).
Die Meldungsart sollte zusätzlich zum internen Kürzel im Klartext erscheinen
(Kriterium~2.4), damit Informationen in vertrautem Vokabular dargestellt
werden,\footcite[Vgl.][Abschnitt~5.2.2.5]{DIN_ISO_9241-110:2020} und die
Oberflächentexte sollten durchgängig deutsch gehalten werden (Kriterium~3.4), um
vertrauten sprachlichen Konventionen zu
entsprechen.\footcite[Vgl.][Abschnitt~5.3.3.1]{DIN_ISO_9241-110:2020}
Beide Maßnahmen senken das Risiko einer Fehlauswahl der Meldung und damit einer
Fehlzuordnung der Fotos.\\
\textbf{Transparenz von Systemzustand und Rückmeldung:}
Mit mittlerer Priorität sollte die Erkennbarkeit des Systemzustands verbessert
werden.
Die aktiven Filter sollten sichtbar und zentral im Suchbereich erscheinen statt
abseits am oberen Fensterrand (Kriterium~2.1), sodass jederzeit erkennbar ist,
welche Einstellungen
wirken.\footcite[Vgl.][Abschnitt~5.2.2.2]{DIN_ISO_9241-110:2020}
Bei länger dauernden Vorgängen wie dem Laden der Meldungsliste oder dem Hochladen
von Fotos sollte eine Warte- oder Fortschrittsanzeige erscheinen (Kriterium~3.3),
damit erkennbar bleibt, dass das System
arbeitet.\footcite[Vgl.][Abschnitt~5.3.2.4]{DIN_ISO_9241-110:2020}\\
\textbf{Konsistenz der Oberflächengestaltung:}
Mit niedriger Priorität empfiehlt sich, die Schaltflächen nach einem einheitlichen
Schema zu gestalten, das primäre und sekundäre Aktionen klar unterscheidet
(Kriterium~2.3), damit ihr Interaktionscharakter unmittelbarer erkennbar
wird.\footcite[Vgl.][Abschnitt~5.2.2.4]{DIN_ISO_9241-110:2020}\\
\textbf{Anpassung an den Einsatzkontext vor Ort:}
Ebenfalls niedrige Priorität hat die Anpassung an die Einsatzbedingungen vor Ort,
insbesondere größere, auch mit Handschuhen sicher bedienbare Bedienelemente
(Kriterium~7.3), um keine unverhältnismäßigen Anforderungen zu
stellen.\footcite[Vgl.][Abschnitt~5.7.2.7]{DIN_ISO_9241-110:2020}

Tabelle~\ref{tab:optimierung} fasst die abgeleiteten Maßnahmen mit
Kriteriumsbezug, Schweregrad und Priorität zusammen.

\begin{table}[htbp]
\centering
\caption{Abgeleitete Optimierungspotenziale im Überblick}
\label{tab:optimierung}
{\footnotesize
\begin{tabularx}{\textwidth}{|c|>{\RaggedRight}X|c|c|}
\hline
\textbf{Krit.} & \textbf{Maßnahme} & \textbf{SG} & \textbf{Priorität} \\\hline
6.5 & Dateiformat prüfen und Upload-Fehler zurückmelden & 3 & hoch \\\hline
6.8 & Fehlertext korrekt anzeigen & 3 & hoch \\\hline
6.9 & Fehlermeldungen mit konkreter Handlungsanweisung versehen & 3 & hoch \\\hline
6.6 & Bestätigungsabfrage vor dem Löschen eines Fotos & 2 & hoch \\\hline
6.3 & Abhängige Filter dynamisch auf gültige Werte einschränken & 2 & hoch \\\hline
3.4 & Oberfläche durchgängig deutschsprachig gestalten & 3 & hoch \\\hline
2.4 & Meldungsart zusätzlich im Klartext anzeigen & 2 & hoch \\\hline
2.1 & Aktive Filter sichtbar und zentral anzeigen & 2 & mittel \\\hline
3.3 & Warte-/Fortschrittsanzeige bei langen Vorgängen & 2 & mittel \\\hline
2.3 & Schaltflächen nach einheitlichem Schema gestalten & 1 & niedrig \\\hline
7.3 & Bedienelemente für Handschuhnutzung vergrößern & 1 & niedrig \\\hline
\end{tabularx}
}
\end{table}
